\section*{Capítulo \ref{ch:b1:structure-spectroscopy} --- Espectroscopia para a estrutura}

\begin{solution}{exo:b1:structure-spectroscopy:1}
$\varepsilon = A/(\ell c) = 0.64/(1.00 \times \num{2.0e-5}) =
\qty{3.2e4}{L/(mol.cm)}$. Com o triplo da concentração: $A = 1.92$, se a
lei ainda valer nessa \omterm{def:b1:structure-spectroscopy:absorbance}{absorbância}.
\end{solution}

\begin{solution}{exo:b1:structure-spectroscopy:2}
\ce{C6H6}: $D = (12 + 2 - 6)/2 = 4$, benzeno (um anel e três ligações $\pi$).
\ce{C4H8O2}: 1, etanoato de etila ou ácido butanoico. \ce{C3H7NO}: $(6 + 2 + 1
- 7)/2 = 1$, propanamida. \ce{C2H3Cl}: 1, cloroeteno. \ce{C10H14O}:
$(20 + 2 - 14)/2 = 4$: a carvona tem um anel, duas C=C e uma C=O.
\end{solution}

\begin{solution}{exo:b1:structure-spectroscopy:3}
Propanona: um sinal (6H). Ácido etanoico: dois (3H, 1H). 2-Metilpropano:
dois (9H, 1H). 1,2-Dicloroetano: um (4H).
\end{solution}

\begin{solution}{exo:b1:structure-spectroscopy:4}
$\delta = 1460/400 = \qty{3.65}{ppm}$; a \qty{90}{MHz}, o sinal fica a
$3.65 \times 90 = \qty{329}{Hz}$ da referência. As \omterm{def:b1:structure-spectroscopy:coupling}{constantes de
acoplamento} permanecem as mesmas em hertz, enquanto as diferenças de
deslocamento crescem com o campo: \omterm{def:b1:structure-spectroscopy:coupling}{multipletos} que se sobrepõem em campo
baixo se separam em campo alto.
\end{solution}

\begin{solution}{exo:b1:structure-spectroscopy:5}
\ce{CH3}: tripleto (3H); \ce{CH2} central: cinco vizinhos, sexteto
1\,:\,5\,:\,10\,:\,10\,:\,5\,:\,1 (2H); \ce{CH2Br}: tripleto (2H), o mais
desblindado.
\end{solution}

\begin{solution}{exo:b1:structure-spectroscopy:6}
Etanol: O--H (3666) e nenhuma C=O; propanona: C=O (1738) e nenhuma O--H;
ácido etanoico: as duas, O--H (3582) e C=O (1788).
\end{solution}

\begin{solution}{exo:b1:structure-spectroscopy:7}
Etanoato de etila: quarteto (2H) perto de 4.1, singleto (3H) perto de
2.0, tripleto (3H) perto de 1.3 ppm. Propanoato de metila: singleto (3H)
do \ce{OCH3}, desblindado pelo oxigênio (perto de 3.7 ppm), quarteto (2H)
do \ce{CH2C=O} (perto de 2.3), tripleto (3H) perto de 1.1. O quarteto é o
sinal mais desblindado no primeiro e não no segundo: um quarteto perto de
4 ppm indica \ce{O-CH2CH3}.
\end{solution}

\begin{solution}{exo:b1:structure-spectroscopy:8}
Linhas separadas por \qty{0.080}{ppm}, ou seja, $0.080 \times 90 = \qty{7.2}{Hz}$:
a \omterm{def:b1:structure-spectroscopy:coupling}{constante de acoplamento}. O tripleto tem o mesmo espaçamento: o
\ce{CH2} e o \ce{CH3} estão acoplados entre si.
\end{solution}

\begin{solution}{exo:b1:structure-spectroscopy:9}
Uma banda C=O e um único sinal: a propanona, \ce{CH3COCH3}. O isômero
aldeído, o propanal \ce{CH3CH2CHO}, mostra o estiramento C--H do grupo
\ce{CHO} e seu próton muito desblindado.
\end{solution}

\begin{solution}{exo:b1:structure-spectroscopy:10}
O primeiro conjunto desdobra a linha em $n_1 + 1$ linhas; cada uma delas
é desdobrada pelo segundo conjunto em $n_2 + 1$: $(n_1 + 1)(n_2 + 1)$
linhas. Se $J_1 =
J_2$, a posição de uma linha só depende de $k_1 + k_2$, que assume
$n_1 + n_2 +
1$ valores: a regra do $n + 1$ com $n = n_1 + n_2$. \ce{CH2} central do
1-bromopropano: $4 \times 3 = 12$ linhas em geral, seis (um sexteto)
quando as duas constantes são iguais.
\end{solution}

\begin{solution}{exo:b1:structure-spectroscopy:11}
$D = 0$; um álcool (banda O--H, nenhuma C=O). Dubleto (6H): dois \ce{CH3}
equivalentes em um mesmo CH; \omterm{def:b1:structure-spectroscopy:coupling}{multipleto} (1H): esse CH; dubleto (2H): um
\ce{CH2} vizinho do CH e que leva o OH; singleto largo: OH.
2-Metilpropan-1-ol, \ce{(CH3)2CH-CH2-OH}.
\end{solution}

\begin{solution}{exo:b1:structure-spectroscopy:12}
Cada espécie absorve independentemente: $A = A_1 + A_2 = \ell(\varepsilon_1c_1
+ \varepsilon_2c_2)$. Em dois comprimentos de onda, com os quatro
coeficientes conhecidos, duas equações lineares dão $c_1$ e $c_2$.
\end{solution}

\begin{solution}{pb:b1:structure-spectroscopy:1}
\textbf{1.} $0.5690 \times 12.0/44.0 = \qty{0.1552}{g}$.
\textbf{2.} $0.2328 \times 2.0/18.0 = \qty{0.02587}{g}$.
\textbf{3.} $0.2500 - 0.1552 - 0.0259 = \qty{0.0689}{g}$.
\textbf{4.} C \qty{62.1}{\percent}, H \qty{10.3}{\percent}, O
\qty{27.6}{\percent}.
\textbf{5.} C $0.1552/12.0 = 0.01293$, H $0.02587$, O $0.0689/16.0 =
0.00431$ mol: razão $3.00 : 6.00 : 1$, \ce{C3H6O}.
\textbf{6.} $M = \rho RT/p = 3740 \times 8.314 \times 373.15/(1.000 \times 10^5) =
\qty{116}{g/mol}$ (massa específica em \unit{g/m^3}).
\textbf{7.} $116/58 = 2$: \ce{C6H12O2}, $D = (12 + 2 - 12)/2 = 1$.
\textbf{8.} Um estiramento C=O.
\textbf{9.} Qualquer O--H: nenhum ácido, nenhum álcool.
\textbf{10.} Um estiramento C--O, forte: o C--O de um éster.
\textbf{11.} Estiramentos C--H dos grupos \ce{CH2} e \ce{CH3}.
\textbf{12.} Uma ligação $\pi$, usada pela C=O; com um C--O forte e
nenhum O--H, um éster. Um ácido exige um O--H; um diol não tem C=O. Uma
cetona com um grupo éter teria sua C=O perto dos \qty{1738}{cm^{-1}} da
propanona, enquanto os \qty{1754}{cm^{-1}} observados ficam perto dos
\qty{1764}{cm^{-1}} do etanoato de etila; a RMN da Parte III confirma o
éster.
\textbf{13.} Cinco; $2 + 2 + 2 + 3 + 3 = 12$, os doze hidrogênios.
\textbf{14.} Um \ce{CH2} com três vizinhos (um \ce{CH3}), ligado ao
oxigênio (desblindado): \ce{-O-CH2-CH3}.
\textbf{15.} O \ce{CH3} desse grupo etila, com dois vizinhos; os dois
sinais estão acoplados entre si, daí um único $J$.
\textbf{16.} Um \ce{CH2} com dois vizinhos, ao lado da C=O.
\textbf{17.} Um \ce{CH2} com cinco vizinhos: entre um \ce{CH2} e um
\ce{CH3}.
\textbf{18.} Um \ce{CH3} com dois vizinhos, longe do grupo funcional.
\textbf{19.} \ce{CH3CH2O-} e \ce{-CH2CH2CH3}.
\textbf{20.} Os prótons em 4.13 ppm estão no carbono ligado ao oxigênio,
os prótons em 2.28 ppm ao lado da C=O: ambos retiradores de elétrons.
\textbf{21.} \ce{CH3CH2-O-CO-CH2CH2CH3} ou \ce{CH3CH2-CO-O-CH2CH2CH3}. Só
o primeiro põe o \ce{CH2} com três vizinhos no oxigênio (o quarteto em
4.13 ppm).
\textbf{22.} O propanoato de propila é o segundo: seu \ce{OCH2} seria um
tripleto e o quarteto ficaria perto de 2.3 ppm.
\textbf{23.} O etanoato de butila, \ce{CH3CO-O-C4H9}, mostraria um
singleto (3H) perto de 2.0 ppm e nenhum quarteto.
\textbf{24.} Butanoato de etila, \ce{CH3CH2CH2-CO-O-CH2CH3}.
\textbf{25.} Infravermelho: C=O 1754, C--O 1182, C--H 2982, nenhum O--H.
RMN: \ce{OCH2} 4.13 (q), \ce{CH2CO} 2.28 (t), \ce{CH2} central 1.66
(sexteto), \ce{CH3} do éster 1.25 (t), \ce{CH3} da cadeia 0.95 (t),
\omterm{def:b1:structure-spectroscopy:integration}{integrações} 2, 2, 2, 3, 3.
\textbf{26.} \textbf{\qty{116}{g/mol}}: butanoato de etila, \ce{C6H12O2}.
\end{solution}
